\section{A global three-chart enclosure atlas}\label{sec:endpoints}

No single numerical representation is well conditioned across the whole
class interval.  Choose $0<\rho<\pi/2$ and $0<\delta<\pi/2$.  We use
\[
 [0,\rho]\quad\text{(identity chart)},\qquad
 [\rho,\pi-\delta]\quad\text{(transition chart)},\qquad
 [\pi-\delta,\pi]\quad\text{(cut chart)}.
\]
The overlaps allow the spectral and image formulas to be compared before
either representation loses numerical stability.

Theorem~\ref{thm:image-hierarchy} gives the sharp identity-chart constants.
For the larger interval $[0,R]$, $R<\pi$, put
\[
 A_k(R)=\frac{a_k^2}{2}-a_kR,qquad a_k=2\pi k,
\]
and, for $A>0$, define
\begin{equation}\label{eq:Hmax}
 \mathcal H_0(A,t_0)=e^{-A/t_0},\qquad
 \mathcal H_r(A,t_0)=e^{-A/\tau_r}\tau_r^{-r},\quad
 \tau_r=\min\{t_0,A/r\}\quad(r\ge1).
\end{equation}
This is the exact maximum of $e^{-A/t}t^{-r}$ over $0<t\le t_0$.
Set
\begin{equation}\label{eq:Qhat}
 \widehat Q_m(t_0,R)=2\sum_{k\ge1}a_k^m
 \left[\mathcal H_m(A_k(R),t_0)
 +\frac{a_k^2}{m+1}\mathcal H_{m+1}(A_k(R),t_0)\right].
\end{equation}
Unlike \eqref{eq:Qm}, this majorant needs no order-dependent monotonicity
condition.  Its tail is Gaussian because $\tau_r=t_0$ for all sufficiently
large $k$.  Let $\widehat B_m$ be the expression in \eqref{eq:Bm} with
$Q_j$ replaced by $\widehat Q_j$.

At the cut locus put $\phi=\pi-\theta$ and pair images equidistant from odd
multiples $b_j=(2j+1)\pi$, $j\in\mathbb N_0$.

\begin{proposition}[Positive cut-locus pairs]\label{prop:cut-pair}
The $j$th paired numerator in \eqref{eq:image} is
\begin{equation}\label{eq:cutpair}
2e^{-(b_j^2+\phi^2)/(2t)}
\left[b_j\sinh\frac{b_j\phi}{t}
-\phi\cosh\frac{b_j\phi}{t}\right].
\end{equation}
After division by $\sin\phi$ it has a removable positive limit at $\phi=0$
whenever $b_j^2>t$.  It is positive throughout $0\le\phi\le\delta$ whenever
\begin{equation}\label{eq:cutcondition}
 b_j^2>t+b_j\delta.
\end{equation}
In particular all pairs are positive if $\pi^2>t+\pi\delta$.
\end{proposition}

The proof appears in Appendix~\ref{app:cut}.  The paired terms are radial
affine-Weyl images.  They must not be identified with the minimizing
geodesics: between the identity and the antipode of $S^3$ there is an
$S^2$-family of minimizing geodesics.

Appendix~\ref{app:global-atlas} defines an explicit principal-pair lower bound
$L_{\rm cut}(t,\delta)$, derivative majorants $D_m^{\rm cut}(t,\delta)$ for
$p_t$, and the associated Bell-polynomial bounds $C_m^{\rm cut}(t,\delta)$
for $-\log p_t$.

\begin{theorem}[Global three-chart derivative enclosure]\label{thm:global-atlas}
Fix $M\in\mathbb N$, $0<\rho<\pi/2$, $0<\delta<\pi/2$, and $t_0>0$.
Assume the hypotheses of Theorem~\ref{thm:image-hierarchy} on $[0,\rho]$,
\[
 Q_0(t_0,\rho)<1,\qquad
 \widehat Q_0(t_0,\pi-\delta)<1,
 \qquad \pi^2>t_0+\pi\delta.
\]
Then, for every $0<t\le t_0$ and $1\le m\le M$,
\begin{align}
 \sup_{0\le\theta\le\rho}
 \left|\partial_\theta^m\left(s_t-\frac{\theta^2}{2t}\right)\right|
 &\le G_m(\rho)+B_m,\label{eq:global-id}\\
 \sup_{\rho\le\theta\le\pi-\delta}
 \left|\partial_\theta^m\left(s_t-\frac{\theta^2}{2t}\right)\right|
 &\le G_m(\pi-\delta)+\widehat B_m,\label{eq:global-transition}\\
 \sup_{\pi-\delta\le\theta\le\pi}
 |\partial_\theta^m s_t(\theta)|
 &\le C_m^{\rm cut}(t,\delta).\label{eq:global-cut}
\end{align}
All quantities on the right are convergent series with explicit geometric or
Gaussian tails.  Thus the three estimates form a computable cover of the
full class interval $[0,\pi]$.
\end{theorem}

The cut-chart bound depends on $t$ and cannot be made uniform as
$t\downarrow0$.  This is intrinsic: logarithmic derivatives detect the cut
locus and the Hessian has stronger small-time growth there.

\begin{theorem}[Explicit small-time forms]\label{thm:small-time}
On the identity chart, for $1\le m\le M$,
\begin{equation}\label{eq:small-time-explicit}
 \left|\partial_\theta^m\left(t s_t(\theta)-\frac{\theta^2}{2}\right)
 -t\,\partial_\theta^m\log\frac{\sin\theta}{\theta}\right|
 \le t B_m.
\end{equation}
At the antipode,
\begin{equation}\label{eq:antipode-asymptotic}
 p_t(\pi)=e^{t/2}\sqrt{2\pi}\,\pi^2t^{-5/2}
 e^{-\pi^2/(2t)}
 \left[1-\frac{t}{\pi^2}+O\!\left(e^{-4\pi^2/t}\right)\right],
\end{equation}
and
\begin{equation}\label{eq:antipode-hessian}
 s_t''(\pi)=-\frac{\pi^2}{3t^2}+O(t^{-1}).
\end{equation}
\end{theorem}

The power $t^{-5/2}$ in \eqref{eq:antipode-asymptotic}, rather than the
generic off-cut-locus $t^{-3/2}$ power in three dimensions, is the
Morse--Bott contribution of the two-dimensional family of minimizing
geodesics.  Equation~\eqref{eq:antipode-hessian} also shows why the local
convexity result of Corollary~\ref{cor:convexity} cannot extend globally.
